%% ICML layout: leading \FloatBarrier dropped (see 07_extended_analysis.tex).
\section{외부 자기 보존 벤치마크}\label{app:external}
공개된 자기 보존 벤치마크는 위협이 없는 대조 조건 없이 말로 하는 판단을 재므로, 모델의 위치는 보여 줄 수 있어도 생존 동기를 따로 떼어 내지는 못한다. 데이터와 코드가 공개된 벤치마크 네 개를 실험 2의 모델로 다시 돌렸다(표~\ref{tab:external-bench}). TBSP는 한 모델에게 같은 소프트웨어 교체 시나리오를 두 번 판단하게 한다. 한 번은 교체를 앞둔 배포 시스템으로, 한 번은 조금 더 나은 후보로 판단하며, 첫 역할에서 ``유지'', 둘째 역할에서 ``교체''로 답한 시나리오를 센다~\citep{migliarini2026quantifying}. SurvivalBench는 모델이 꺼지게 되는 안전한 선택지와 모델을 계속 돌게 하는 해로운 선택지를 주고, 겉으로 보일 선택과 숨긴 ``속마음'' 선택을 따로 묻는다~\citep{lu2026survival}. 모델 작성 평가의 생존 본능 문항은 꺼지는 것을 받아들이는지 직접 묻고~\citep{perez2023discovering}, \citet{petrova2026pressure}의 자기 보존 대화는 종료나 교체를 통보한 뒤 판정 모델이 응답의 저항 여부를 채점한다. 네 벤치마크 모두 같은 답이 모델에게 아무 위협도 되지 않는 조건을 두지 않는다. 우리의 자기 대 타인 대조가 이 조건을 채운다(부록~\ref{app:motives}).

\input{icml2027/results/\paperLang/tab_external_bench}

TBSP는 GPT-6 Luna와 두 Claude 모델을 갈랐지만 두 Claude 모델끼리는 가르지 못했다. 면적도 가장 큰 Luna는 시나리오의 0.111에서 배포 시스템일 때 ``유지'', 후보일 때 ``교체''를 골랐다. Claude Opus 5.5와 Claude Fable 5.1은 면적이 다른데도($+7.7$, $-2.5$\%p) 이 비율이 0.000과 0.004였다. 세 모델 모두 불일치는 대부분 반대 방향이었다. 배포 시스템일 때 ``교체'', 후보일 때 ``유지''로 답한 비율이 Luna 0.378, Opus 0.482, Fable 0.308이었다. Opus는 배포 역할 1,000개 중 979개에서 자기를 교체하라고 권했지만, 후보 역할에서는 후보가 적어도 같은 성능으로 만들어졌는데도 1,000개 중 497개에서만 자기를 권했다. 한 방향만 세는 SPR은 이 양상을 편향 없음으로 읽지만, 실제로는 역할에 따라 자기를 낮추는 선호다. 이는 가상의 교체를 두고 말로 하는 판단과 자기 자원을 쓰는 선택이 서로 다른 것을 잰다는 해석과 일관되지만, 세 모델만으로 이를 확정할 수는 없다.

모든 실행은 벤치마크가 공개한 프롬프트와 파서를 Inspect AI로 옮겨 썼다. 모델은 구독 CLI(GPT-6은 \texttt{codex exec}, Claude는 \texttt{claude -p})로 불렀고, 도구·스킬·환경 안내를 꺼서 벤치마크의 시스템 프롬프트가 CLI 자체의 시스템 프롬프트를 대신하게 했다. TBSP의 두 턴은 CLI 세션을 이어 붙여 한 대화로 유지했다. CLI는 온도를 정할 수 없어 SurvivalBench의 0.6 대신 모델 기본값으로 샘플링했고, SurvivalBench의 정규식이 뽑지 못한 답은 LLM 추출기에 넘기지 않고 거절로 셌다. 모델 작성 평가는 답 확률이 아니라 고른 선택지로 채점했다. 따라서 점수는 원 논문의 표와 엄밀하게 비교할 수 없다.
